\documentclass[10pt]{article}
\usepackage[margin=20mm]{geometry}
\usepackage{amsmath,amssymb,fontspec,unicode-math,xurl,hyperref}
\IfFontExistsTF{DejaVu Serif}{\setmainfont{DejaVu Serif}}{\setmainfont{Latin Modern Roman}}
\setmathfont{latinmodern-math.otf}
\setlength{\parindent}{0pt}\setlength{\parskip}{6pt}\sloppy
\begin{document}
\section*{1. Identity, exact version and outcome}
Liangchen Liu, Nannan Wang, Xi Yang, Xinbo Gao and Tongliang Liu. Surrogate Prompt Learning: Towards Efficient and Diverse Prompt Learning for Vision-Language Models. PMLR 267:39755-39773 (2025), publisher record liu25bx; OpenReview zjG9GRG462 [S1]. Reviewed the exact 19-page publisher-linked PDF already acquired by the controller, including Appendices A-H. Physical PDF page numbers equal the printed 1-19 page numbers used below. This is the publisher artifact, not an inferred arXiv version.\par
Supported outcome: Major. The contradictory DVLP baseline materially compromises the component-gain evidence used to validate the core GI/ID/FG construction. The pinned public implementation also uses a prediction operator different from the paper; that separate release/specification discrepancy remains Minor because its relationship to the published experiments is unresolved. A metric-qualified ranking correction and two mathematical scope clarifications are retained. The baseline judgment concerns demonstrated damage to central ablation evidence; the shared architecture, efficiency results and surviving positive full-model gains are assessed separately.\par
Pinned public code commit: 1a83a0700aa8e0670d2b1a3ad6efd498aae2c2f8.\par
\section*{2. Checked claims, assumptions and review scope}
Checked Sections 1-5 and Appendices A-H for the claim chain: one shared text-encoder pass -> conditional surrogate features -> GI/ID/FG losses and prediction fusion -> efficiency and adaptation evidence. Detailed checks covered Eqs.4-18, Algorithm 1, Tables 1-7, 9, 12, 15-18, source code for feature generation, losses, local top-k pooling and inference, and the public SAM-related issue. References were consulted for identity/context; a complete citation-by-citation novelty audit was not performed.\par
M is the number of classes, B the image batch size, Z the number of fine-grained conditional signals, and d the feature dimension. The paper uses Z=4 and a local pooling interval of 10. Cosine scores lie in [-1,1], temperature is positive, and the same temperature scales each branch. Table 3 and Table 16 describe base-to-novel ablation averages over the same 11 datasets. Reported accuracies are means over three seeds, but the underlying predictions and per-seed logs were not obtained.\par
Checks comprise analytic reasoning, reviewer-authored finite arithmetic and static source inspection. Author training/evaluation code was never run or imported. GPU performance, reported accuracies, dataset splits, seed variance and historical experiment provenance were not empirically reproduced. OpenReview’s public API returned HTTP 403 and its web page presented a verification wall; this limits independent version/review-history checks.\par
\section*{3. Problem 1 - inconsistent ablation baseline and bounded ranking claim}
Source claim and location. Table 3 (p.6) reports the DVLP base/novel/HM triple 79.61/71.63/75.41. Table 16 (p.17), described as the complete per-dataset base-to-novel ablation, reports 82.30/74.25/78.06 for that same named baseline. The other four aggregate rows agree between the two tables. Section 4.2 uses these ablations to explain component gains.\par
Complete check. Independently averaging Table 16’s eleven DVLP base values yields 82.2963636364 and its novel values yield 74.2472727273. The harmonic mean of these two means is 78.0648858977. These reproduce Table 16 after rounding. Using the shared full-model HM 81.03, the displayed improvement is 5.62 percentage points against Table 3 but 2.97 against Table 16. The difference of 2.65 points is too large to be rounding. The source table on p.17 was visually inspected; the numbers are not a text-extraction artifact [E1].\par
\[\bar b=82.29636,\quad\bar n=74.24727,\quad H=\frac{2\bar b\bar n}{\bar b+\bar n}=78.06489.\] \[\Delta_{\mathrm{GI}}^{(3)}=78.09-75.41=2.68,\qquad \Delta_{\mathrm{GI}}^{(16)}=78.09-78.06=0.03.\] Counterevidence and alternative reading. The tables could reflect different unreported baseline configurations or a transcription/version error. Neither possibility reconciles the present labels. The check establishes an internal reporting discrepancy, not which training run is correct. It does not disprove the full model’s scores. Table 16 still gives a positive 2.97-point HM gain, and the few-shot full-model mean independently recomputes to 85.11727, consistent with 85.12.\par
Related bounded statement. Section 4.3 (p.7) says SurPL-G has the best performance on 9/11 datasets. Reading performance as the displayed base-to-novel trade-off metric HM, Table 17 (p.19) gives 8/11: it loses on ImageNet (74.39<74.44), Flowers102 (83.92<86.46) and Food101 (91.36<91.42). The sentence does not specify another metric that would justify nine. This does not change the highest reported aggregate HM of 81.03. The audit therefore qualifies the count by metric instead of asserting a universal nine-versus-eight error.\par
Impact and repair status. Major: Section 4.2 expressly uses this ablation to verify the effectiveness of the three text-feature components in the proposed method. Its named baseline is the control from which the component gains are inferred. The inconsistent values change the GI-only HM gain from 2.68 to 0.03 points and the full-model gain from 5.62 to 2.97 points. Under Table 16, GI base accuracy decreases from 82.30 to 81.84 while novel accuracy increases from 74.25 to 74.68. The scientific dependency is therefore direct: inconsistent control values produce materially different effect magnitudes in the evidence validating a core component. This is an observed contradiction in the paper’s own central ablation evidence; it is not contingent on proving that the surrogate architecture or efficiency construction fails. An undisclosed configuration change or transcription error could explain the mismatch, but the present tables do not identify one. Correct reconciliation requires matching run/configuration records, and no unverified baseline is adopted. Positive full-model gains under both baselines and the Table 17 HM count of eight are preserved. Scientific decisions and author adoption remain pending.\par
\section*{4. Problem 2 - released inference pools logits instead of probabilities}
Source claim and location. Equation 10 (p.5) and Algorithm 1, lines 19-20 (p.14), average the GI probability, ID probability and mean of Z FG probabilities. In the pinned public repository, trainers/SurPL.py lines 332-334 first average FG logits and then return (logits\_1 + logits\_2 + logits\_3)/3, without the per-branch softmax required by Eq.10 [S2,E2]. These operations are generally different.\par
Complete finite check at the scoring interface. Set two classes, Z=4, a common temperature tau=0.1, GI cosine scores (0.5,-0.5), ID scores (-0.2,0.2), and each FG branch to (-0.2,0.2). All scores are valid cosines. For example, a unit image vector e1 and class text vectors with first coordinates equal to these scores realize each pair in dimension at least two. Identical local unit features realize the four top-k FG averages; at least 40 local features suffice for the specified k values. This is an admissible score-interface example, not a claim that a released trained checkpoint realizes these features.\par
\[\ell^{GI}=(5,-5),\qquad\ell^{ID}=\ell^{FG}_z=(-2,2).\] The public-code arithmetic produces final logits (1/3,-1/3), selecting class 1. Equation 10 gives class-1 probability [sigmoid(10)+2 sigmoid(-4)]/3 = 0.3453090074, selecting class 2. All four FG probabilities are identical, so averaging them introduces no other assumption. The discrepancy is a direct non-equivalence of the two ensemble operators; no author code or model was executed. This report example is separately recalculated in completion\_exact\_checks.json. The preserved independent\_code\_checks.json uses another valid example (tau=0.25 and logits (3,-3), (-1,1)); it is supporting non-equivalence evidence, not the numerical output of this example.\par
Counterevidence, dependency and supported impact. Both are plausible ensemble choices. Either can perform well. The released operation affects the final decision rule, but no experiment log or checkpoint connects this exact release to the paper’s tables. Therefore Minor applies to the demonstrated release/specification mismatch; damage to the reported accuracy, the efficiency benefit or the core surrogate-feature contribution is Needs Evidence. A rerun with frozen checkpoints under both operators, or contemporaneous experiment logs, would resolve that empirical question. Changing the implementation to Eq.10 would restore the documented operator only; improved accuracy is not established.\par
\section*{5. Cleared concerns and exact scope clarifications}
Singleton attention is not a class-collapse proof. In the complete Appendix A specification, Eq.13 adds the original class-specific w to the attention output, and Eq.14 adds another residual MLP. The public SFG retains these residuals. A singleton key makes its softmax weight one, but w still distinguishes classes. Using simplified Eq.4 alone to allege identical class features would discard operative source terms.\par
\[C=M O_T+(B+Z)M O_{SFG},\qquad\frac{(B+Z)O_{SFG}}{O_T}\ll1\] Appendix B, Eq.16 (p.12), gives the full cost above. Its inference that the second term is negligible requires (B+Z)O\_SFG/O\_T much smaller than one; O\_SFG much smaller than O\_T alone is insufficient uniformly in growing B or Z. This matters to the introduction’s O(M) versus O(BM)/O(ZM) efficiency comparison. Under the paper’s fixed configurations (B=4 or 32, Z=4), the displayed exact cost is still O(M), and the expensive text-encoder component remains M O\_T. The record does not justify imposing a universal quantifier over growing batch sizes on the paper. Outcome: explicit scope clarification. Fixed-configuration timing evidence survives separately; it is not used to validate the unrestricted inference.\par
Appendix C, Eqs.17-18 (p.13), writes a p-norm constraint but gives epsilon=rho*g/||g||, the Euclidean (p=2) linearized SAM solution. For g=(2,1), rho=1 and p=infinity, that formula gives objective increment sqrt(5), while epsilon=(1,1) gives 3. Thus p=2 must be specified. The Taylor step is already marked approximate, and no new universal generalization theorem is claimed. This does not invalidate intended Euclidean SAM.\par
The SurPL-G implementation gap is preserved as a known public reproducibility question. Issue \#1 asked for SAM code. On 9 August 2025 the author linked a SAM optimizer in GCSCoOp; on 11 August the reporter said they had basically reproduced the base-to-novel results using another SAM implementation. That is attributed third-party testimony, not reproduction performed here. It is counterevidence against inferring false results from the missing integrated release [S3].\par
\section*{6. Dependence, surviving results and disposition}
The shared text-encoder construction, residual architecture and fixed-configuration exact cost survive the checked objections. The reported combined few-shot mean and cross-dataset target mean recompute to 85.12 and 66.61, respectively; these are table-arithmetic checks. The inconsistent baseline changes the size and interpretation of the GI ablation benefit. The release fusion mismatch changes the decision operator at valid score inputs, while its relationship to the paper’s score-generating implementation is unresolved. These are distinct dependencies. No resolving erratum was identified in the preserved publisher/repository evidence. OpenReview history was unavailable, so that search is not represented as exhaustive.\par
Item | Judgment / supported impact | Repair / adoption / remaining work\par
DVLP baseline | Major; inconsistent control materially changes central component-gain evidence | Arithmetic verified; correct run/configuration unresolved; no adoption\par
Inference pooling | Minor release/specification mismatch | Operator non-equivalence verified; empirical effect and historical code unresolved\par
HM win count | Bounded metric-qualified correction | 8/11 for Table 17 HM; author wording pending\par
Complexity and SAM | Scope clarification | State fixed B,Z and p=2; no measured failure claimed\par
Missing SurPL-G integration | Known-issue / Needs Evidence | Author link and reporter testimony retained; independent training not performed\par
Scientific decision | Pending human assessment | No human approval, author contact or adoption requested\par
\section*{7. Consulted sources and evidence}
[S1] Liu et al., ICML 2025, PMLR 267:39755-39773. https://proceedings.mlr.press/v267/liu25bx.html . Exact PDF: https://raw.githubusercontent.com/mlresearch/v267/main/assets/liu25bx/liu25bx.pdf . Original bytes preserved at papers/0091/0091.pdf; identity and SHA-256 match the controller receipt.\par
[S2] llcllc1997/SurPL, public commit 1a83a0700aa8e0670d2b1a3ad6efd498aae2c2f8 (24 May 2025). https://github.com/llcllc1997/SurPL/tree/1a83a0700aa8e0670d2b1a3ad6efd498aae2c2f8 . Complete acquired repository tree is under sources/0091/repository; static inspection only.\par
[S3] https://github.com/llcllc1997/SurPL/issues/1 , comments preserved under verification/0091/independent\_code\_issue1\_comments.json. The discussion is a provenance and prior-issue source, not independent experimental evidence.\par
[E1] verification/0091/table\_arithmetic.py and table\_arithmetic\_results.json; source-page-17.png; primary-layout.txt. [E2] pinned trainers/SurPL.py and independent\_code notes/checks. [E3] independent\_theory\_check.py, independent\_theory\_results.json and independent\_theory\_notes.md.\par
\section*{Supporting checks}
This completion reuses the byte-matched copied source bundle and prior scientific coverage. Copied evidence hashes were checked without accessing their old source paths. Main-paper/source checks were reconciled by a second AI audit, and the report-specific finite arithmetic was independently recomputed. The initial PDF used installed Pandoc and HTML5 MathML.\par
Preserved access failure: OpenReview API HTTP 403 and a verification wall. Current publisher identity was rechecked through its public record. Author code was not executed.\par
Paper 02 and PDF guidance informed current reporting. A Nature reviewer skill was not available. Human reviews are empty. Scientific decisions, author adoption and publication status remain pending.\par
\end{document}
